% 04-resimler.tex

% Resim sayfaları için özel kenar boşlukları ayarlanır
\newgeometry{top=2.5cm, left=2.5cm, right=1.5cm, bottom=1cm}

% Resim sayfalarında sayfa numaralandırması "Mevcut Sayfa / Toplam Sayfa" şeklinde olmalıdır
\fancyhf{} 
\renewcommand{\headrulewidth}{0pt}
\cfoot{\thepage \ / \pageref{ResimlerSonSayfa}}

\begin{center}
    % ÖNEMLİ BİLGİ: 
    % Türkpatent kurallarına göre resimlerde (siyah dışında) renk ve gölgelendirme kullanılmaz. 
    % Şekillerin içi boyalı veya siyah dolu olmamalıdır (tarama çizgileri kullanılabilir).
    % Resimlerde antet veya sayfa çerçevesi bulunmamalıdır.
    
    % Şekil 1
    % Resim dosyalarınızı 'sekiller/' klasörüne yükleyip adını aşağıya yazınız.
    \includegraphics[width=0.8\textwidth]{sekil1} % Uzantı (.png, .jpg) yazmanıza gerek yoktur
    \vspace{2em}
    
    % Eğer aynı sayfaya birden fazla şekil ekleyecekseniz:
    % \includegraphics[width=0.8\textwidth]{sekil2}
    
\end{center}

% Yeni bir resim sayfasına geçmek için \newpage komutunu kullanabilirsiniz:
% \newpage
% \begin{center}
%     \includegraphics[width=0.8\textwidth]{sekil3}
% \end{center}

% Bu etiket, toplam resim sayfası sayısını hesaplamak için en sona eklenmiştir. Lütfen silmeyiniz.
\label{ResimlerSonSayfa}
